%! Quadratic Systems \& Modeling — Unit 2, Lesson 2.5 — Student Packet Cover Sheet
% SLIDES-FIRST shape (2026-09-25), condensed Unit 2 (2026-09-29). 12pt; \coverbanner measures the
% title block. TWO scored rows: the slide handout (the student's notes, the pages right after this
% cover) and the homework — a DeltaMath set ONLINE, not in the packet (user choice).
% Homework is due the first class after two study halls — never "next class".
\documentclass[12pt]{article}
\usepackage{algebra2-article}
\usepackage{algebra2-boxes}

\begin{document}

\coverbanner{Unit 2 \quad Quadratic Functions}{Lesson 2.5 \quad Quadratic Systems \& Modeling}

\namedateperiod
\vspace{0.18in}

\boxguard[14]
\begin{learningtargetbox}
\small
\begin{enumerate}[leftmargin=1.5em, itemsep=5pt]
  \item \ldots{} solve a \textbf{linear--quadratic system} by \textbf{substitution}, giving each
        \textbf{solution} as an ordered pair --- an \textbf{intersection point} of the graphs.
  \item \ldots{} predict whether a line and a parabola meet in $0$, $1$, or $2$ points from the
        \textbf{discriminant} of the collapsed equation.
  \item \ldots{} answer a modeling question with the right feature --- the
        \textbf{$y$-intercept}, a \textbf{zero}, or the \textbf{vertex} --- and tell the
        \textbf{maximum value} (an output) from \emph{when} it happens (an input).
\end{enumerate}
\end{learningtargetbox}

\vspace{0.14in}

\boxguard[14]
\begin{tocbox}
\small
\renewcommand{\arraystretch}{1.5}
\begin{tabularx}{\linewidth}{@{} c l X r @{}}
\rowcolor{forestbg}
\textbf{\#} & \textbf{Component} & \textbf{Description} & \textbf{Score} \\
\midrule
1 & Slide Notes & The printed slides: the warm-up, three rounds of definitions, examples, and
                  practice ladders, and the Final Round --- worked in the notes column & \blank{1.2cm} \\
2 & DeltaMath   & \textbf{A2: Lesson 2.5: Systems \& Modeling} --- online, not in this packet;
                  \textbf{scored; due the first class after two study halls}      & \blank{1.2cm} \\
\midrule
  & \multicolumn{2}{r}{\textbf{Total}} & \blank{1.2cm} \\
\end{tabularx}
\end{tocbox}

\vspace{0.14in}

\boxguard[8]
\begin{remindbox}
\small
A \textbf{solution of a system} is an ordered pair that makes \emph{both} equations true --- a
point where the graphs meet. Solve by \textbf{substitution}: set the $y$'s equal, solve the one
quadratic, then put each $x$ back into the line for its $y$. A line and a parabola meet in $2$,
$1$, or $0$ points, as the collapsed equation's \textbf{discriminant} is positive, zero, or
negative. In a model, the $y$-intercept is the starting value, the positive zero is when it
reaches $0$, and the vertex is the best value: the \textbf{maximum or minimum value} is the
vertex's \emph{output} (how high, how much); \emph{when} or \emph{where} it happens is the input.
\end{remindbox}

\end{document}
